%% ====================================================================
%%  Metodología de evaluación
%% ====================================================================
\section{Metodología de evaluación}\label{sec:metodologia}

La metodología de evaluación aplicada para la elección de las
tecnologías de proxy-caché y almacenamiento de archivos de gran tamaño
fue DAR (Decision Analysis and Resolution) del modelo CMMI (Capability
Maturity Model Integration) (CMMI Institute, 2010). Este referente
metodológico permitió evaluar las necesidades, problemas y oportunidades
(NPO) del grupo de investigación GRID a través de un proceso
estructurado que consideró múltiples alternativas, criterios de
evaluación bien definidos y un análisis comparativo cuantitativo. En
este caso, se identificaron y analizaron 14 tecnologías proxy-caché
como: HAProxy, NGINX, Envoy, Varnish, Traefik, Caddy, Apache Traffic
Server, Kong, Squid, Istio, Linkerd, K8Sidecar, Proxy-Service y nuster y
22 tecnologías de almacenamiento incluyendo Ceph, Lustre, JuiceFS, DAOS,
SeaweedFS, HDFS, MinIO, Longhorn, entre otras; aplicando criterios
ponderados como rendimiento técnico (25 \%), escalabilidad (20 \%),
soporte y mantenimiento (20 \%), adopción y madurez (10 \%),
observabilidad (10 \%), adaptación al GRID y riesgo tecnológico (5 \%).
Mediante el uso de matrices de evaluación NPO y puntuación DAR, se logró
visualizar las fortalezas y debilidades de cada opción, facilitando una
selección alineada con los objetivos estratégicos del GRID. Así, el uso
de DAR no solo aportó transparencia al proceso, sino también
trazabilidad y justificación técnica frente a una decisión crítica para
la arquitectura de infraestructura del grupo de investigación.
